\documentclass[10pt,a4paper]{article}
\usepackage[margin=2.2cm]{geometry}
\usepackage{mathptmx}
\usepackage{xcolor}
\usepackage{enumitem}
\usepackage{amsmath}
\usepackage{parskip}
\usepackage{xr}
\usepackage{url}
\externaldocument{access}
\newcommand{\AblMainCostDeg}{5.0}
\newcommand{\AblMainCostF}{78.8}
\newcommand{\AblMainCostRest}{0.0}
\newcommand{\AblMainCostTerm}{6.3}
\newcommand{\AblMainCostTot}{135.5}
\newcommand{\AblMainCostV}{45.4}
\newcommand{\AblNoDegConvMax}{5.3}
\newcommand{\AblNoDegCost}{145.6}
\newcommand{\AblNoDegCostDeg}{5.0}
\newcommand{\AblNoDegCostF}{80.2}
\newcommand{\AblNoDegCostRest}{0.1}
\newcommand{\AblNoDegCostSd}{10.3}
\newcommand{\AblNoDegCostTerm}{13.9}
\newcommand{\AblNoDegCostTot}{145.6}
\newcommand{\AblNoDegCostV}{46.3}
\newcommand{\AblNoDegN}{3}
\newcommand{\AblNoDegNadirMeanHz}{47.8}
\newcommand{\AblNoDegNadirMeanHzSd}{0.3}
\newcommand{\AblNoDegOvPuH}{0.002}
\newcommand{\AblNoDegOvPuHSd}{0.001}
\newcommand{\AblNoDegPairDiffs}{$-$1.0, $-$2.2, $-$2.0}
\newcommand{\AblNoDegPairMean}{-1.7}
\newcommand{\AblNoDegPval}{0.85}
\newcommand{\AblNoDegReserveMeanMw}{23.1}
\newcommand{\AblNoDegReserveMeanMwSd}{1.9}
\newcommand{\AblNoDegReserveRedPct}{-1}
\newcommand{\AblNoDegSeedCosts}{139.5, 157.4, 139.8}
\newcommand{\AblNoDegSeedSocDefs}{6.9, 26.3, 8.4}
\newcommand{\AblNoDegSocDefMaxSeed}{26.3}
\newcommand{\AblNoDegSocDefPct}{13.88}
\newcommand{\AblNoDegSocDefPctSd}{10.75}
\newcommand{\AblNoDegThroughputMwh}{452}
\newcommand{\AblNoDegThroughputMwhSd}{23}
\newcommand{\AblNoDegUvPuH}{0.92}
\newcommand{\AblNoDegUvPuHSd}{0.03}
\newcommand{\AblNoFConvMax}{0.8}
\newcommand{\AblNoFCost}{169.0}
\newcommand{\AblNoFCostDeg}{5.7}
\newcommand{\AblNoFCostF}{89.6}
\newcommand{\AblNoFCostRest}{2.2}
\newcommand{\AblNoFCostSd}{0.9}
\newcommand{\AblNoFCostTerm}{28.7}
\newcommand{\AblNoFCostTot}{169.0}
\newcommand{\AblNoFCostV}{42.9}
\newcommand{\AblNoFN}{3}
\newcommand{\AblNoFNadirMeanHz}{51.0}
\newcommand{\AblNoFNadirMeanHzSd}{0.4}
\newcommand{\AblNoFOvPuH}{0.001}
\newcommand{\AblNoFOvPuHSd}{0.000}
\newcommand{\AblNoFPairDiffs}{+29.4, +8.4, +27.3}
\newcommand{\AblNoFPairMean}{21.7}
\newcommand{\AblNoFPval}{0.07}
\newcommand{\AblNoFReserveMeanMw}{10.1}
\newcommand{\AblNoFReserveMeanMwSd}{1.3}
\newcommand{\AblNoFReserveRedPct}{56}
\newcommand{\AblNoFSeedCosts}{170.0, 168.1, 169.1}
\newcommand{\AblNoFSeedSocDefs}{28.7, 29.0, 28.3}
\newcommand{\AblNoFSocDefMaxSeed}{29.0}
\newcommand{\AblNoFSocDefPct}{28.68}
\newcommand{\AblNoFSocDefPctSd}{0.36}
\newcommand{\AblNoFThroughputMwh}{512}
\newcommand{\AblNoFThroughputMwhSd}{29}
\newcommand{\AblNoFUvPuH}{0.86}
\newcommand{\AblNoFUvPuHSd}{0.01}
\newcommand{\AblNoOVConvMax}{11.4}
\newcommand{\AblNoOVCost}{144.4}
\newcommand{\AblNoOVCostDeg}{4.8}
\newcommand{\AblNoOVCostF}{79.4}
\newcommand{\AblNoOVCostRest}{0.0}
\newcommand{\AblNoOVCostSd}{5.4}
\newcommand{\AblNoOVCostTerm}{13.8}
\newcommand{\AblNoOVCostTot}{144.4}
\newcommand{\AblNoOVCostV}{46.4}
\newcommand{\AblNoOVN}{3}
\newcommand{\AblNoOVNadirMeanHz}{47.5}
\newcommand{\AblNoOVNadirMeanHzSd}{0.2}
\newcommand{\AblNoOVOvPuH}{0.004}
\newcommand{\AblNoOVOvPuHSd}{0.001}
\newcommand{\AblNoOVPairDiffs}{+0.5, $-$9.1, $-$0.3}
\newcommand{\AblNoOVPairMean}{-2.9}
\newcommand{\AblNoOVPval}{0.70}
\newcommand{\AblNoOVReserveMeanMw}{24.4}
\newcommand{\AblNoOVReserveMeanMwSd}{0.5}
\newcommand{\AblNoOVReserveRedPct}{-6}
\newcommand{\AblNoOVSeedCosts}{141.1, 150.6, 141.5}
\newcommand{\AblNoOVSeedSocDefs}{8.6, 21.8, 10.9}
\newcommand{\AblNoOVSocDefMaxSeed}{21.8}
\newcommand{\AblNoOVSocDefPct}{13.75}
\newcommand{\AblNoOVSocDefPctSd}{7.06}
\newcommand{\AblNoOVThroughputMwh}{431}
\newcommand{\AblNoOVThroughputMwhSd}{32}
\newcommand{\AblNoOVUvPuH}{0.92}
\newcommand{\AblNoOVUvPuHSd}{0.02}
\newcommand{\AblRefConvMax}{10.8}
\newcommand{\AblRefCost}{147.3}
\newcommand{\AblRefCostDeg}{4.9}
\newcommand{\AblRefCostF}{80.3}
\newcommand{\AblRefCostRest}{0.2}
\newcommand{\AblRefCostSd}{10.7}
\newcommand{\AblRefCostTerm}{15.6}
\newcommand{\AblRefCostTot}{147.3}
\newcommand{\AblRefCostV}{46.4}
\newcommand{\AblRefN}{3}
\newcommand{\AblRefNadirMeanHz}{47.8}
\newcommand{\AblRefNadirMeanHzSd}{0.4}
\newcommand{\AblRefOvPuH}{0.003}
\newcommand{\AblRefOvPuHSd}{0.002}
\newcommand{\AblRefReserveMeanMw}{22.9}
\newcommand{\AblRefReserveMeanMwSd}{2.6}
\newcommand{\AblRefReserveRedPct}{0}
\newcommand{\AblRefSeedCosts}{140.5, 159.7, 141.8}
\newcommand{\AblRefSeedSocDefs}{7.7, 27.8, 11.4}
\newcommand{\AblRefSocDefMaxSeed}{27.8}
\newcommand{\AblRefSocDefPct}{15.62}
\newcommand{\AblRefSocDefPctSd}{10.68}
\newcommand{\AblRefThroughputMwh}{437}
\newcommand{\AblRefThroughputMwhSd}{25}
\newcommand{\AblRefUvPuH}{0.92}
\newcommand{\AblRefUvPuHSd}{0.03}
\newcommand{\AblSquashConvMax}{0.4}
\newcommand{\AblSquashCost}{178.4}
\newcommand{\AblSquashCostDeg}{5.0}
\newcommand{\AblSquashCostF}{83.6}
\newcommand{\AblSquashCostRest}{2.0}
\newcommand{\AblSquashCostSd}{0.4}
\newcommand{\AblSquashCostTerm}{28.1}
\newcommand{\AblSquashCostTot}{178.4}
\newcommand{\AblSquashCostV}{59.6}
\newcommand{\AblSquashN}{3}
\newcommand{\AblSquashNadirMeanHz}{48.9}
\newcommand{\AblSquashNadirMeanHzSd}{0.1}
\newcommand{\AblSquashOvPuH}{0.010}
\newcommand{\AblSquashOvPuHSd}{0.000}
\newcommand{\AblSquashPairDiffs}{+38.2, +18.7, +36.2}
\newcommand{\AblSquashPairMean}{31.0}
\newcommand{\AblSquashPval}{0.04}
\newcommand{\AblSquashReserveMeanMw}{18.6}
\newcommand{\AblSquashReserveMeanMwSd}{0.2}
\newcommand{\AblSquashReserveRedPct}{19}
\newcommand{\AblSquashSeedCosts}{178.7, 178.4, 178.0}
\newcommand{\AblSquashSeedSocDefs}{27.6, 28.6, 28.2}
\newcommand{\AblSquashSocDefMaxSeed}{28.6}
\newcommand{\AblSquashSocDefPct}{28.14}
\newcommand{\AblSquashSocDefPctSd}{0.48}
\newcommand{\AblSquashThroughputMwh}{454}
\newcommand{\AblSquashThroughputMwhSd}{18}
\newcommand{\AblSquashUvPuH}{1.18}
\newcommand{\AblSquashUvPuHSd}{0.01}
\newcommand{\AblWfHighConvMax}{4.3}
\newcommand{\AblWfHighCost}{140.5}
\newcommand{\AblWfHighCostDeg}{4.6}
\newcommand{\AblWfHighCostF}{77.9}
\newcommand{\AblWfHighCostRest}{0.0}
\newcommand{\AblWfHighCostSd}{0.4}
\newcommand{\AblWfHighCostTerm}{7.8}
\newcommand{\AblWfHighCostTot}{140.5}
\newcommand{\AblWfHighCostV}{50.2}
\newcommand{\AblWfHighN}{2}
\newcommand{\AblWfHighNadirMeanHz}{46.9}
\newcommand{\AblWfHighNadirMeanHzSd}{0.4}
\newcommand{\AblWfHighOvPuH}{0.007}
\newcommand{\AblWfHighOvPuHSd}{0.001}
\newcommand{\AblWfHighPairDiffs}{$-$0.3, $-$18.9}
\newcommand{\AblWfHighPairMean}{-9.6}
\newcommand{\AblWfHighPval}{0.38}
\newcommand{\AblWfHighReserveMeanMw}{26.6}
\newcommand{\AblWfHighReserveMeanMwSd}{1.3}
\newcommand{\AblWfHighReserveRedPct}{-16}
\newcommand{\AblWfHighSeedCosts}{140.2, 140.8}
\newcommand{\AblWfHighSeedSocDefs}{7.2, 8.5}
\newcommand{\AblWfHighSocDefMaxSeed}{8.5}
\newcommand{\AblWfHighSocDefPct}{7.84}
\newcommand{\AblWfHighSocDefPctSd}{0.90}
\newcommand{\AblWfHighThroughputMwh}{415}
\newcommand{\AblWfHighThroughputMwhSd}{21}
\newcommand{\AblWfHighUvPuH}{1.00}
\newcommand{\AblWfHighUvPuHSd}{0.01}
\newcommand{\AblWfLowConvMax}{4.8}
\newcommand{\AblWfLowCost}{161.9}
\newcommand{\AblWfLowCostDeg}{5.1}
\newcommand{\AblWfLowCostF}{84.7}
\newcommand{\AblWfLowCostRest}{1.1}
\newcommand{\AblWfLowCostSd}{5.8}
\newcommand{\AblWfLowCostTerm}{27.4}
\newcommand{\AblWfLowCostTot}{161.9}
\newcommand{\AblWfLowCostV}{43.6}
\newcommand{\AblWfLowN}{2}
\newcommand{\AblWfLowNadirMeanHz}{49.3}
\newcommand{\AblWfLowNadirMeanHzSd}{0.8}
\newcommand{\AblWfLowOvPuH}{0.001}
\newcommand{\AblWfLowOvPuHSd}{0.001}
\newcommand{\AblWfLowPairDiffs}{+17.3, +6.3}
\newcommand{\AblWfLowPairMean}{11.8}
\newcommand{\AblWfLowPval}{0.14}
\newcommand{\AblWfLowReserveMeanMw}{16.5}
\newcommand{\AblWfLowReserveMeanMwSd}{3.4}
\newcommand{\AblWfLowReserveRedPct}{28}
\newcommand{\AblWfLowSeedCosts}{157.9, 166.0}
\newcommand{\AblWfLowSeedSocDefs}{25.6, 29.2}
\newcommand{\AblWfLowSocDefMaxSeed}{29.2}
\newcommand{\AblWfLowSocDefPct}{27.38}
\newcommand{\AblWfLowSocDefPctSd}{2.56}
\newcommand{\AblWfLowThroughputMwh}{458}
\newcommand{\AblWfLowThroughputMwhSd}{30}
\newcommand{\AblWfLowUvPuH}{0.87}
\newcommand{\AblWfLowUvPuHSd}{0.00}
\newcommand{\AblWvHighConvMax}{1.7}
\newcommand{\AblWvHighCost}{167.5}
\newcommand{\AblWvHighCostDeg}{4.9}
\newcommand{\AblWvHighCostF}{86.6}
\newcommand{\AblWvHighCostRest}{2.3}
\newcommand{\AblWvHighCostSd}{0.4}
\newcommand{\AblWvHighCostTerm}{30.8}
\newcommand{\AblWvHighCostTot}{167.5}
\newcommand{\AblWvHighCostV}{42.9}
\newcommand{\AblWvHighN}{2}
\newcommand{\AblWvHighNadirMeanHz}{49.9}
\newcommand{\AblWvHighNadirMeanHzSd}{0.1}
\newcommand{\AblWvHighOvPuH}{0.003}
\newcommand{\AblWvHighOvPuHSd}{0.002}
\newcommand{\AblWvHighPairDiffs}{+27.2, +7.6}
\newcommand{\AblWvHighPairMean}{17.4}
\newcommand{\AblWvHighPval}{0.08}
\newcommand{\AblWvHighReserveMeanMw}{13.9}
\newcommand{\AblWvHighReserveMeanMwSd}{0.3}
\newcommand{\AblWvHighReserveRedPct}{39}
\newcommand{\AblWvHighSeedCosts}{167.8, 167.2}
\newcommand{\AblWvHighSeedSocDefs}{31.0, 30.6}
\newcommand{\AblWvHighSocDefMaxSeed}{31.0}
\newcommand{\AblWvHighSocDefPct}{30.85}
\newcommand{\AblWvHighSocDefPctSd}{0.28}
\newcommand{\AblWvHighThroughputMwh}{442}
\newcommand{\AblWvHighThroughputMwhSd}{19}
\newcommand{\AblWvHighUvPuH}{0.85}
\newcommand{\AblWvHighUvPuHSd}{0.01}
\newcommand{\AblWvLowConvMax}{1.8}
\newcommand{\AblWvLowCost}{140.1}
\newcommand{\AblWvLowCostDeg}{4.8}
\newcommand{\AblWvLowCostF}{78.1}
\newcommand{\AblWvLowCostRest}{0.0}
\newcommand{\AblWvLowCostSd}{0.4}
\newcommand{\AblWvLowCostTerm}{2.2}
\newcommand{\AblWvLowCostTot}{140.1}
\newcommand{\AblWvLowCostV}{55.1}
\newcommand{\AblWvLowN}{2}
\newcommand{\AblWvLowNadirMeanHz}{47.0}
\newcommand{\AblWvLowNadirMeanHzSd}{0.1}
\newcommand{\AblWvLowOvPuH}{0.008}
\newcommand{\AblWvLowOvPuHSd}{0.003}
\newcommand{\AblWvLowPairDiffs}{$-$0.7, $-$19.3}
\newcommand{\AblWvLowPairMean}{-10.0}
\newcommand{\AblWvLowPval}{0.36}
\newcommand{\AblWvLowReserveMeanMw}{26.4}
\newcommand{\AblWvLowReserveMeanMwSd}{0.4}
\newcommand{\AblWvLowReserveRedPct}{-15}
\newcommand{\AblWvLowSeedCosts}{139.9, 140.4}
\newcommand{\AblWvLowSeedSocDefs}{2.7, 1.6}
\newcommand{\AblWvLowSocDefMaxSeed}{2.7}
\newcommand{\AblWvLowSocDefPct}{2.17}
\newcommand{\AblWvLowSocDefPctSd}{0.78}
\newcommand{\AblWvLowThroughputMwh}{428}
\newcommand{\AblWvLowThroughputMwhSd}{3}
\newcommand{\AblWvLowUvPuH}{1.09}
\newcommand{\AblWvLowUvPuHSd}{0.03}
\newcommand{\BestValCost}{137.5}
\newcommand{\ClipAbovePsetOne}{13.4}
\newcommand{\ClipAbovePsetTwo}{21.6}
\newcommand{\ClipAboveRhoOne}{43.2}
\newcommand{\ClipAboveRhoTwo}{30.7}
\newcommand{\ClipBelowPsetOne}{25.1}
\newcommand{\ClipBelowPsetTwo}{33.9}
\newcommand{\ClipBelowRhoOne}{10.8}
\newcommand{\ClipBelowRhoTwo}{17.9}
\newcommand{\ClipDimAll}{49.2}
\newcommand{\ClipDimPsetOne}{38.5}
\newcommand{\ClipDimPsetTwo}{55.5}
\newcommand{\ClipDimRhoOne}{54.0}
\newcommand{\ClipDimRhoTwo}{48.6}
\newcommand{\ClipDimSteps}{3\,600}
\newcommand{\ClipMedianExcess}{0.61}
\newcommand{\ClipSquashAll}{0.0}
\newcommand{\ConvChangePct}{1.0}
\newcommand{\CrossDAOHi}{264\,192}
\newcommand{\CrossDAOLo}{104\,448}
\newcommand{\DayPPOPeakOne}{19}
\newcommand{\DayRBNightVmin}{0.91}
\newcommand{\DayResEveMin}{13.9}
\newcommand{\DayResHi}{33}
\newcommand{\DayResLo}{23}
\newcommand{\DayUCVminOne}{0.84}
\newcommand{\FinalValCost}{137.7}
\newcommand{\FinalValCostSd}{1.6}
\newcommand{\FreqHour}{13:00}
\newcommand{\HrHtwo}{3.25}
\newcommand{\HrPPONadir}{74.4}
\newcommand{\HrPPONadirOne}{74.4}
\newcommand{\HrPPONadirTwo}{47.7}
\newcommand{\HrPPORes}{32.3}
\newcommand{\HrRBNadir}{73.5}
\newcommand{\HrRBNadirOne}{73.5}
\newcommand{\HrRBNadirTwo}{57.0}
\newcommand{\HrRBRes}{19.1}
\newcommand{\HrUCNadir}{78.6}
\newcommand{\HrUCNadirOne}{78.6}
\newcommand{\HrUCNadirTwo}{76.9}
\newcommand{\HrUCRes}{0.0}
\newcommand{\ModeDecomposedDAOCost}{152.3}
\newcommand{\ModeDecomposedDAOEViolMwh}{0.000}
\newcommand{\ModeDecomposedDAONadirMeanHz}{49.8}
\newcommand{\ModeDecomposedDAORegBattMwh}{-0.43}
\newcommand{\ModeDecomposedDAOSocDefPct}{9.68}
\newcommand{\ModeDecomposedDAOUnmetMwh}{1.28}
\newcommand{\ModeDecomposedDAOUvPuH}{1.045}
\newcommand{\ModeDecomposedDAOViolHours}{18.2}
\newcommand{\ModeDecomposedPPOCost}{135.5}
\newcommand{\ModeDecomposedPPOEViolMwh}{0.000}
\newcommand{\ModeDecomposedPPONadirMeanHz}{47.3}
\newcommand{\ModeDecomposedPPORegBattMwh}{-0.69}
\newcommand{\ModeDecomposedPPOSocDefPct}{6.31}
\newcommand{\ModeDecomposedPPOUnmetMwh}{0.00}
\newcommand{\ModeDecomposedPPOUvPuH}{0.904}
\newcommand{\ModeDecomposedPPOViolHours}{11.9}
\newcommand{\ModeDecomposedRBCost}{245.5}
\newcommand{\ModeDecomposedRBEViolMwh}{0.000}
\newcommand{\ModeDecomposedRBNadirMeanHz}{49.4}
\newcommand{\ModeDecomposedRBRegBattMwh}{-0.46}
\newcommand{\ModeDecomposedRBSocDefPct}{13.58}
\newcommand{\ModeDecomposedRBUnmetMwh}{0.00}
\newcommand{\ModeDecomposedRBUvPuH}{2.860}
\newcommand{\ModeDecomposedRBViolHours}{12.5}
\newcommand{\ModeDecomposedUCCost}{477.2}
\newcommand{\ModeDecomposedUCEViolMwh}{0.000}
\newcommand{\ModeDecomposedUCNadirMeanHz}{54.4}
\newcommand{\ModeDecomposedUCRegBattMwh}{0.00}
\newcommand{\ModeDecomposedUCSocDefPct}{0.00}
\newcommand{\ModeDecomposedUCUnmetMwh}{0.00}
\newcommand{\ModeDecomposedUCUvPuH}{7.402}
\newcommand{\ModeDecomposedUCViolHours}{19.3}
\newcommand{\ModeDeltaCostDecomposed}{0.0}
\newcommand{\ModeDeltaCostFullhour}{-0.0}
\newcommand{\ModeDeltaCostHour}{0.1}
\newcommand{\ModeDeltaCostNone}{-0.1}
\newcommand{\ModeDeltaCostSdDecomposed}{0.0}
\newcommand{\ModeDeltaCostSdFullhour}{-0.0}
\newcommand{\ModeDeltaCostSdHour}{0.1}
\newcommand{\ModeDeltaCostSdNone}{-0.0}
\newcommand{\ModeDeltaCostSdWindow}{-0.0}
\newcommand{\ModeDeltaCostWindow}{-0.1}
\newcommand{\ModeFullhourDAOCost}{152.3}
\newcommand{\ModeFullhourDAOEViolMwh}{0.000}
\newcommand{\ModeFullhourDAONadirMeanHz}{49.8}
\newcommand{\ModeFullhourDAORegBattMwh}{-0.47}
\newcommand{\ModeFullhourDAOSocDefPct}{9.68}
\newcommand{\ModeFullhourDAOUnmetMwh}{1.28}
\newcommand{\ModeFullhourDAOUvPuH}{1.045}
\newcommand{\ModeFullhourDAOViolHours}{18.2}
\newcommand{\ModeFullhourPPOCost}{135.5}
\newcommand{\ModeFullhourPPOEViolMwh}{0.000}
\newcommand{\ModeFullhourPPONadirMeanHz}{47.3}
\newcommand{\ModeFullhourPPORegBattMwh}{-0.73}
\newcommand{\ModeFullhourPPOSocDefPct}{6.30}
\newcommand{\ModeFullhourPPOUnmetMwh}{0.00}
\newcommand{\ModeFullhourPPOUvPuH}{0.904}
\newcommand{\ModeFullhourPPOViolHours}{11.9}
\newcommand{\ModeFullhourRBCost}{245.1}
\newcommand{\ModeFullhourRBEViolMwh}{0.000}
\newcommand{\ModeFullhourRBNadirMeanHz}{49.4}
\newcommand{\ModeFullhourRBRegBattMwh}{-0.50}
\newcommand{\ModeFullhourRBSocDefPct}{13.75}
\newcommand{\ModeFullhourRBUnmetMwh}{0.00}
\newcommand{\ModeFullhourRBUvPuH}{2.848}
\newcommand{\ModeFullhourRBViolHours}{12.5}
\newcommand{\ModeFullhourUCCost}{477.2}
\newcommand{\ModeFullhourUCEViolMwh}{0.000}
\newcommand{\ModeFullhourUCNadirMeanHz}{54.4}
\newcommand{\ModeFullhourUCRegBattMwh}{0.00}
\newcommand{\ModeFullhourUCSocDefPct}{0.00}
\newcommand{\ModeFullhourUCUnmetMwh}{0.00}
\newcommand{\ModeFullhourUCUvPuH}{7.402}
\newcommand{\ModeFullhourUCViolHours}{19.3}
\newcommand{\ModeHourDAOCost}{152.5}
\newcommand{\ModeHourDAOEViolMwh}{0.000}
\newcommand{\ModeHourDAONadirMeanHz}{49.8}
\newcommand{\ModeHourDAORegBattMwh}{-1.02}
\newcommand{\ModeHourDAOSocDefPct}{9.92}
\newcommand{\ModeHourDAOUnmetMwh}{1.30}
\newcommand{\ModeHourDAOUvPuH}{1.045}
\newcommand{\ModeHourDAOViolHours}{18.2}
\newcommand{\ModeHourPPOCost}{135.6}
\newcommand{\ModeHourPPOEViolMwh}{0.000}
\newcommand{\ModeHourPPONadirMeanHz}{47.3}
\newcommand{\ModeHourPPORegBattMwh}{-1.67}
\newcommand{\ModeHourPPOSocDefPct}{6.42}
\newcommand{\ModeHourPPOUnmetMwh}{0.00}
\newcommand{\ModeHourPPOUvPuH}{0.905}
\newcommand{\ModeHourPPOViolHours}{12.0}
\newcommand{\ModeHourRBCost}{247.0}
\newcommand{\ModeHourRBEViolMwh}{0.000}
\newcommand{\ModeHourRBNadirMeanHz}{49.4}
\newcommand{\ModeHourRBRegBattMwh}{-1.10}
\newcommand{\ModeHourRBSocDefPct}{13.66}
\newcommand{\ModeHourRBUnmetMwh}{0.00}
\newcommand{\ModeHourRBUvPuH}{2.887}
\newcommand{\ModeHourRBViolHours}{12.6}
\newcommand{\ModeHourUCCost}{477.2}
\newcommand{\ModeHourUCEViolMwh}{0.000}
\newcommand{\ModeHourUCNadirMeanHz}{54.4}
\newcommand{\ModeHourUCRegBattMwh}{0.00}
\newcommand{\ModeHourUCSocDefPct}{0.00}
\newcommand{\ModeHourUCUnmetMwh}{0.00}
\newcommand{\ModeHourUCUvPuH}{7.402}
\newcommand{\ModeHourUCViolHours}{19.3}
\newcommand{\ModeNoneDAOCost}{152.2}
\newcommand{\ModeNoneDAOEViolMwh}{0.000}
\newcommand{\ModeNoneDAONadirMeanHz}{49.8}
\newcommand{\ModeNoneDAORegBattMwh}{0.00}
\newcommand{\ModeNoneDAOSocDefPct}{9.59}
\newcommand{\ModeNoneDAOUnmetMwh}{1.28}
\newcommand{\ModeNoneDAOUvPuH}{1.045}
\newcommand{\ModeNoneDAOViolHours}{18.2}
\newcommand{\ModeNonePPOCost}{135.5}
\newcommand{\ModeNonePPOEViolMwh}{0.000}
\newcommand{\ModeNonePPONadirMeanHz}{47.3}
\newcommand{\ModeNonePPORegBattMwh}{0.00}
\newcommand{\ModeNonePPOSocDefPct}{6.26}
\newcommand{\ModeNonePPOUnmetMwh}{0.00}
\newcommand{\ModeNonePPOUvPuH}{0.904}
\newcommand{\ModeNonePPOViolHours}{11.9}
\newcommand{\ModeNoneRBCost}{245.0}
\newcommand{\ModeNoneRBEViolMwh}{0.000}
\newcommand{\ModeNoneRBNadirMeanHz}{49.4}
\newcommand{\ModeNoneRBRegBattMwh}{0.00}
\newcommand{\ModeNoneRBSocDefPct}{13.71}
\newcommand{\ModeNoneRBUnmetMwh}{0.00}
\newcommand{\ModeNoneRBUvPuH}{2.848}
\newcommand{\ModeNoneRBViolHours}{12.5}
\newcommand{\ModeNoneUCCost}{477.2}
\newcommand{\ModeNoneUCEViolMwh}{0.000}
\newcommand{\ModeNoneUCNadirMeanHz}{54.4}
\newcommand{\ModeNoneUCRegBattMwh}{0.00}
\newcommand{\ModeNoneUCSocDefPct}{0.00}
\newcommand{\ModeNoneUCUnmetMwh}{0.00}
\newcommand{\ModeNoneUCUvPuH}{7.402}
\newcommand{\ModeNoneUCViolHours}{19.3}
\newcommand{\ModeWindowDAOCost}{152.2}
\newcommand{\ModeWindowDAOEViolMwh}{0.000}
\newcommand{\ModeWindowDAONadirMeanHz}{49.8}
\newcommand{\ModeWindowDAORegBattMwh}{-0.02}
\newcommand{\ModeWindowDAOSocDefPct}{9.59}
\newcommand{\ModeWindowDAOUnmetMwh}{1.28}
\newcommand{\ModeWindowDAOUvPuH}{1.045}
\newcommand{\ModeWindowDAOViolHours}{18.2}
\newcommand{\ModeWindowPPOCost}{135.5}
\newcommand{\ModeWindowPPOEViolMwh}{0.000}
\newcommand{\ModeWindowPPONadirMeanHz}{47.3}
\newcommand{\ModeWindowPPORegBattMwh}{-0.03}
\newcommand{\ModeWindowPPOSocDefPct}{6.26}
\newcommand{\ModeWindowPPOUnmetMwh}{0.00}
\newcommand{\ModeWindowPPOUvPuH}{0.904}
\newcommand{\ModeWindowPPOViolHours}{11.9}
\newcommand{\ModeWindowRBCost}{245.0}
\newcommand{\ModeWindowRBEViolMwh}{0.000}
\newcommand{\ModeWindowRBNadirMeanHz}{49.4}
\newcommand{\ModeWindowRBRegBattMwh}{-0.02}
\newcommand{\ModeWindowRBSocDefPct}{13.72}
\newcommand{\ModeWindowRBUnmetMwh}{0.00}
\newcommand{\ModeWindowRBUvPuH}{2.848}
\newcommand{\ModeWindowRBViolHours}{12.5}
\newcommand{\ModeWindowUCCost}{477.2}
\newcommand{\ModeWindowUCEViolMwh}{0.000}
\newcommand{\ModeWindowUCNadirMeanHz}{54.4}
\newcommand{\ModeWindowUCRegBattMwh}{0.00}
\newcommand{\ModeWindowUCSocDefPct}{0.00}
\newcommand{\ModeWindowUCUnmetMwh}{0.00}
\newcommand{\ModeWindowUCUvPuH}{7.402}
\newcommand{\ModeWindowUCViolHours}{19.3}
\newcommand{\NSeed}{5}
\newcommand{\NStress}{20}
\newcommand{\NTest}{30}
\newcommand{\PPOSeedSd}{1.3}
\newcommand{\RawClipComponents}{14400}
\newcommand{\RawClipFracPct}{49.2}
\newcommand{\RawClipPPct}{47.0}
\newcommand{\RawClipRhoPct}{51.3}
\newcommand{\RegClosureMax}{3.6e-15}
\newcommand{\RegDAOCumMax}{0.55}
\newcommand{\RegDAOCumMean}{-0.43}
\newcommand{\RegDAOCumMin}{-1.40}
\newcommand{\RegDAOGrossDay}{1.3}
\newcommand{\RegDAOHrMax}{0.398}
\newcommand{\RegDAOHrMean}{0.056}
\newcommand{\RegDAOPosDaysPct}{13}
\newcommand{\RegDAOShare}{0.15}
\newcommand{\RegDAOShareGross}{0.40}
\newcommand{\RegDAOSocAreaMax}{0.21}
\newcommand{\RegDAOSocPct}{0.16}
\newcommand{\RegDAOWinCum}{-0.0172}
\newcommand{\RegPPOCumMax}{0.89}
\newcommand{\RegPPOCumMean}{-0.69}
\newcommand{\RegPPOCumMin}{-2.22}
\newcommand{\RegPPOGrossDay}{2.2}
\newcommand{\RegPPOHrMax}{0.607}
\newcommand{\RegPPOHrMean}{0.091}
\newcommand{\RegPPOPosDaysPct}{13}
\newcommand{\RegPPOShare}{0.18}
\newcommand{\RegPPOShareGross}{0.51}
\newcommand{\RegPPOSocAreaMax}{0.38}
\newcommand{\RegPPOSocPct}{0.25}
\newcommand{\RegPPOWinCum}{-0.0278}
\newcommand{\RegRBCumMax}{0.55}
\newcommand{\RegRBCumMean}{-0.46}
\newcommand{\RegRBCumMin}{-1.40}
\newcommand{\RegRBGrossDay}{1.4}
\newcommand{\RegRBHrMax}{0.398}
\newcommand{\RegRBHrMean}{0.059}
\newcommand{\RegRBPosDaysPct}{10}
\newcommand{\RegRBShare}{0.16}
\newcommand{\RegRBShareGross}{0.43}
\newcommand{\RegRBSocAreaMax}{0.21}
\newcommand{\RegRBSocPct}{0.16}
\newcommand{\RegRBWinCum}{-0.0184}
\newcommand{\RegUCCumMax}{0.00}
\newcommand{\RegUCCumMean}{0.00}
\newcommand{\RegUCCumMin}{0.00}
\newcommand{\RegUCGrossDay}{0.0}
\newcommand{\RegUCHrMax}{0.000}
\newcommand{\RegUCHrMean}{0.000}
\newcommand{\RegUCPosDaysPct}{0}
\newcommand{\RegUCShare}{0.00}
\newcommand{\RegUCShareGross}{0.00}
\newcommand{\RegUCSocAreaMax}{0.00}
\newcommand{\RegUCSocPct}{0.00}
\newcommand{\RegUCWinCum}{0.0000}
\newcommand{\RhoEveOne}{0.46}
\newcommand{\RhoMidOne}{0.11}
\newcommand{\RhoMidTwo}{0.58}
\newcommand{\RhoNightOne}{0.60}
\newcommand{\RhoNightTwo}{0.42}
\newcommand{\SDAOCost}{308.3}
\newcommand{\SDAOCostSd}{19.6}
\newcommand{\SDAOEViol}{0.000}
\newcommand{\SDAOFRmsMhz}{17.23}
\newcommand{\SDAOFRmsMhzSd}{1.35}
\newcommand{\SDAOLossMwh}{4.398}
\newcommand{\SDAOLossMwhSd}{0.152}
\newcommand{\SDAONadirMaxHz}{108.6}
\newcommand{\SDAONadirMaxHzSd}{10.0}
\newcommand{\SDAONadirMeanHz}{77.6}
\newcommand{\SDAONadirMeanHzSd}{6.3}
\newcommand{\SDAOOvPuH}{0.000}
\newcommand{\SDAOOvPuHSd}{0.001}
\newcommand{\SDAOProj}{56}
\newcommand{\SDAOProjEvents}{2.8}
\newcommand{\SDAOProjEventsSd}{3.0}
\newcommand{\SDAOPwExcess}{0.000}
\newcommand{\SDAORawClip}{0}
\newcommand{\SDAOReserveMeanMw}{8.8}
\newcommand{\SDAOReserveMeanMwSd}{1.8}
\newcommand{\SDAORocofMax}{0.215}
\newcommand{\SDAORocofMaxSd}{0.018}
\newcommand{\SDAOSocDefPct}{24.30}
\newcommand{\SDAOSocDefPctSd}{8.70}
\newcommand{\SDAOThroughputMwh}{294}
\newcommand{\SDAOThroughputMwhSd}{26}
\newcommand{\SDAOTieRmsMw}{3.57}
\newcommand{\SDAOTieRmsMwSd}{0.26}
\newcommand{\SDAOUnmetMax}{7.4}
\newcommand{\SDAOUnmetMean}{2.43}
\newcommand{\SDAOUnmetMwh}{2.43}
\newcommand{\SDAOUnmetMwhSd}{2.76}
\newcommand{\SDAOUvPuH}{2.88}
\newcommand{\SDAOUvPuHSd}{0.29}
\newcommand{\SDAOViolHours}{18.6}
\newcommand{\SDAOViolHoursSd}{4.5}
\newcommand{\SDAOVmaxDay}{1.048}
\newcommand{\SDAOVmaxDaySd}{0.002}
\newcommand{\SDAOVminDay}{0.903}
\newcommand{\SDAOVminDaySd}{0.004}
\newcommand{\SPDAO}{1.9\times10^{-6}}
\newcommand{\SPPOCost}{283.8}
\newcommand{\SPPOCostSd}{16.2}
\newcommand{\SPPOEViol}{0.000}
\newcommand{\SPPOFRmsMhz}{16.58}
\newcommand{\SPPOFRmsMhzSd}{1.24}
\newcommand{\SPPOLossMwh}{4.393}
\newcommand{\SPPOLossMwhSd}{0.143}
\newcommand{\SPPONadirMaxHz}{105.8}
\newcommand{\SPPONadirMaxHzSd}{9.3}
\newcommand{\SPPONadirMeanHz}{75.0}
\newcommand{\SPPONadirMeanHzSd}{6.1}
\newcommand{\SPPOOvPuH}{0.006}
\newcommand{\SPPOOvPuHSd}{0.009}
\newcommand{\SPPOProj}{0}
\newcommand{\SPPOProjEvents}{0.0}
\newcommand{\SPPOProjEventsSd}{0.0}
\newcommand{\SPPOPwExcess}{0.000}
\newcommand{\SPPORawClip}{4380}
\newcommand{\SPPOReserveMeanMw}{15.3}
\newcommand{\SPPOReserveMeanMwSd}{1.1}
\newcommand{\SPPORocofMax}{0.205}
\newcommand{\SPPORocofMaxSd}{0.019}
\newcommand{\SPPOSocDefPct}{18.02}
\newcommand{\SPPOSocDefPctSd}{4.07}
\newcommand{\SPPOThroughputMwh}{352}
\newcommand{\SPPOThroughputMwhSd}{17}
\newcommand{\SPPOTieRmsMw}{3.57}
\newcommand{\SPPOTieRmsMwSd}{0.26}
\newcommand{\SPPOUnmetMax}{0.0}
\newcommand{\SPPOUnmetMean}{0.00}
\newcommand{\SPPOUnmetMwh}{0.00}
\newcommand{\SPPOUnmetMwhSd}{0.00}
\newcommand{\SPPOUvPuH}{2.66}
\newcommand{\SPPOUvPuHSd}{0.27}
\newcommand{\SPPOViolHours}{18.5}
\newcommand{\SPPOViolHoursSd}{3.0}
\newcommand{\SPPOVmaxDay}{1.052}
\newcommand{\SPPOVmaxDaySd}{0.005}
\newcommand{\SPPOVminDay}{0.905}
\newcommand{\SPPOVminDaySd}{0.004}
\newcommand{\SPRB}{1.9\times10^{-6}}
\newcommand{\SPUC}{1.9\times10^{-6}}
\newcommand{\SRBCost}{369.4}
\newcommand{\SRBCostSd}{17.6}
\newcommand{\SRBEViol}{0.000}
\newcommand{\SRBFRmsMhz}{17.02}
\newcommand{\SRBFRmsMhzSd}{1.30}
\newcommand{\SRBLossMwh}{4.504}
\newcommand{\SRBLossMwhSd}{0.176}
\newcommand{\SRBNadirMaxHz}{108.6}
\newcommand{\SRBNadirMaxHzSd}{9.9}
\newcommand{\SRBNadirMeanHz}{76.8}
\newcommand{\SRBNadirMeanHzSd}{6.2}
\newcommand{\SRBOvPuH}{0.000}
\newcommand{\SRBOvPuHSd}{0.001}
\newcommand{\SRBProj}{0}
\newcommand{\SRBProjEvents}{0.0}
\newcommand{\SRBProjEventsSd}{0.0}
\newcommand{\SRBPwExcess}{0.000}
\newcommand{\SRBRawClip}{0}
\newcommand{\SRBReserveMeanMw}{10.6}
\newcommand{\SRBReserveMeanMwSd}{0.5}
\newcommand{\SRBRocofMax}{0.215}
\newcommand{\SRBRocofMaxSd}{0.018}
\newcommand{\SRBSocDefPct}{26.16}
\newcommand{\SRBSocDefPctSd}{3.65}
\newcommand{\SRBThroughputMwh}{191}
\newcommand{\SRBThroughputMwhSd}{13}
\newcommand{\SRBTieRmsMw}{3.57}
\newcommand{\SRBTieRmsMwSd}{0.26}
\newcommand{\SRBUnmetMax}{0.0}
\newcommand{\SRBUnmetMean}{0.00}
\newcommand{\SRBUnmetMwh}{0.00}
\newcommand{\SRBUnmetMwhSd}{0.00}
\newcommand{\SRBUvPuH}{4.15}
\newcommand{\SRBUvPuHSd}{0.36}
\newcommand{\SRBViolHours}{14.2}
\newcommand{\SRBViolHoursSd}{1.3}
\newcommand{\SRBVmaxDay}{1.043}
\newcommand{\SRBVmaxDaySd}{0.005}
\newcommand{\SRBVminDay}{0.898}
\newcommand{\SRBVminDaySd}{0.003}
\newcommand{\SRedCostDAO}{8.0}
\newcommand{\SRedCostRB}{23.2}
\newcommand{\SRedCostUC}{57.7}
\newcommand{\SRedNadirMaxHzDAO}{2.6}
\newcommand{\SRedNadirMaxHzRB}{2.6}
\newcommand{\SRedNadirMaxHzUC}{12.8}
\newcommand{\SRedNadirMeanHzDAO}{3.4}
\newcommand{\SRedNadirMeanHzRB}{2.4}
\newcommand{\SRedNadirMeanHzUC}{9.2}
\newcommand{\SRedRocofMaxDAO}{5.0}
\newcommand{\SRedRocofMaxRB}{5.0}
\newcommand{\SRedRocofMaxUC}{14.8}
\newcommand{\SRedThroughputMwhDAO}{-20.1}
\newcommand{\SRedThroughputMwhRB}{-84.8}
\newcommand{\SRedThroughputMwhUC}{28.5}
\newcommand{\SRedUvPuHDAO}{7.7}
\newcommand{\SRedUvPuHRB}{35.8}
\newcommand{\SRedUvPuHUC}{74.1}
\newcommand{\SRedViolHoursDAO}{0.6}
\newcommand{\SRedViolHoursRB}{-29.7}
\newcommand{\SRedViolHoursUC}{24.0}
\newcommand{\SUCCost}{671.4}
\newcommand{\SUCCostSd}{24.3}
\newcommand{\SUCEViol}{0.000}
\newcommand{\SUCFRmsMhz}{18.35}
\newcommand{\SUCFRmsMhzSd}{1.38}
\newcommand{\SUCLossMwh}{6.400}
\newcommand{\SUCLossMwhSd}{0.241}
\newcommand{\SUCNadirMaxHz}{121.3}
\newcommand{\SUCNadirMaxHzSd}{11.5}
\newcommand{\SUCNadirMeanHz}{82.6}
\newcommand{\SUCNadirMeanHzSd}{6.4}
\newcommand{\SUCOvPuH}{0.070}
\newcommand{\SUCOvPuHSd}{0.024}
\newcommand{\SUCProj}{329}
\newcommand{\SUCProjEvents}{16.4}
\newcommand{\SUCProjEventsSd}{1.0}
\newcommand{\SUCPwExcess}{0.000}
\newcommand{\SUCRawClip}{0}
\newcommand{\SUCReserveMeanMw}{0.0}
\newcommand{\SUCReserveMeanMwSd}{0.0}
\newcommand{\SUCRocofMax}{0.240}
\newcommand{\SUCRocofMaxSd}{0.021}
\newcommand{\SUCSocDefPct}{0.00}
\newcommand{\SUCSocDefPctSd}{0.00}
\newcommand{\SUCThroughputMwh}{493}
\newcommand{\SUCThroughputMwhSd}{18}
\newcommand{\SUCTieRmsMw}{3.60}
\newcommand{\SUCTieRmsMwSd}{0.25}
\newcommand{\SUCUnmetMax}{0.0}
\newcommand{\SUCUnmetMean}{0.00}
\newcommand{\SUCUnmetMwh}{0.00}
\newcommand{\SUCUnmetMwhSd}{0.00}
\newcommand{\SUCUvPuH}{10.25}
\newcommand{\SUCUvPuHSd}{0.50}
\newcommand{\SUCViolHours}{24.3}
\newcommand{\SUCViolHoursSd}{2.3}
\newcommand{\SUCVmaxDay}{1.070}
\newcommand{\SUCVmaxDaySd}{0.004}
\newcommand{\SUCVminDay}{0.835}
\newcommand{\SUCVminDaySd}{0.006}
\newcommand{\SWinsDAO}{20}
\newcommand{\SWinsRB}{20}
\newcommand{\SWinsUC}{20}
\newcommand{\TDAOCost}{152.3}
\newcommand{\TDAOCostSd}{31.8}
\newcommand{\TDAOEViol}{0.000}
\newcommand{\TDAOFRmsMhz}{11.57}
\newcommand{\TDAOFRmsMhzSd}{1.18}
\newcommand{\TDAOLossMwh}{3.217}
\newcommand{\TDAOLossMwhSd}{0.338}
\newcommand{\TDAONadirMaxHz}{70.5}
\newcommand{\TDAONadirMaxHzSd}{7.6}
\newcommand{\TDAONadirMeanHz}{49.8}
\newcommand{\TDAONadirMeanHzSd}{6.0}
\newcommand{\TDAOOvPuH}{0.001}
\newcommand{\TDAOOvPuHSd}{0.005}
\newcommand{\TDAOProj}{48}
\newcommand{\TDAOProjEvents}{1.6}
\newcommand{\TDAOProjEventsSd}{2.8}
\newcommand{\TDAOPwExcess}{0.000}
\newcommand{\TDAORawClip}{0}
\newcommand{\TDAOReserveMeanMw}{13.9}
\newcommand{\TDAOReserveMeanMwSd}{2.6}
\newcommand{\TDAORocofMax}{0.137}
\newcommand{\TDAORocofMaxSd}{0.016}
\newcommand{\TDAOSocDefMedian}{0.21}
\newcommand{\TDAOSocDefPct}{9.68}
\newcommand{\TDAOSocDefPctSd}{13.55}
\newcommand{\TDAOSocDefZeroPct}{3}
\newcommand{\TDAOThroughputMwh}{348}
\newcommand{\TDAOThroughputMwhSd}{26}
\newcommand{\TDAOTieRmsMw}{2.49}
\newcommand{\TDAOTieRmsMwSd}{0.24}
\newcommand{\TDAOUnmetMax}{10.1}
\newcommand{\TDAOUnmetMean}{1.28}
\newcommand{\TDAOUnmetMwh}{1.28}
\newcommand{\TDAOUnmetMwhSd}{2.55}
\newcommand{\TDAOUvPuH}{1.04}
\newcommand{\TDAOUvPuHSd}{0.37}
\newcommand{\TDAOViolHours}{18.2}
\newcommand{\TDAOViolHoursSd}{4.0}
\newcommand{\TDAOVmaxDay}{1.029}
\newcommand{\TDAOVmaxDaySd}{0.019}
\newcommand{\TDAOVminDay}{0.922}
\newcommand{\TDAOVminDaySd}{0.007}
\newcommand{\TNadirPDAO}{1.9\times10^{-9}}
\newcommand{\TNadirPRB}{1.9\times10^{-9}}
\newcommand{\TNadirPUC}{1.9\times10^{-9}}
\newcommand{\TNadirWinsDAO}{30}
\newcommand{\TNadirWinsRB}{30}
\newcommand{\TNadirWinsUC}{30}
\newcommand{\TPDAO}{1.9\times10^{-9}}
\newcommand{\TPPOCost}{135.5}
\newcommand{\TPPOCostSd}{26.0}
\newcommand{\TPPOEViol}{0.000}
\newcommand{\TPPOFRmsMhz}{10.87}
\newcommand{\TPPOFRmsMhzSd}{1.10}
\newcommand{\TPPOLossMwh}{3.159}
\newcommand{\TPPOLossMwhSd}{0.357}
\newcommand{\TPPONadirMaxHz}{69.6}
\newcommand{\TPPONadirMaxHzSd}{7.5}
\newcommand{\TPPONadirMeanHz}{47.3}
\newcommand{\TPPONadirMeanHzSd}{5.9}
\newcommand{\TPPOOvPuH}{0.004}
\newcommand{\TPPOOvPuHSd}{0.016}
\newcommand{\TPPOProj}{0}
\newcommand{\TPPOProjEvents}{0.0}
\newcommand{\TPPOProjEventsSd}{0.0}
\newcommand{\TPPOPwExcess}{0.000}
\newcommand{\TPPORawClip}{7080}
\newcommand{\TPPOReserveMeanMw}{25.3}
\newcommand{\TPPOReserveMeanMwSd}{2.8}
\newcommand{\TPPORocofMax}{0.133}
\newcommand{\TPPORocofMaxSd}{0.015}
\newcommand{\TPPOSocDefMedian}{4.23}
\newcommand{\TPPOSocDefPct}{6.31}
\newcommand{\TPPOSocDefPctSd}{6.80}
\newcommand{\TPPOSocDefZeroPct}{33}
\newcommand{\TPPOThroughputMwh}{446}
\newcommand{\TPPOThroughputMwhSd}{39}
\newcommand{\TPPOTieRmsMw}{2.51}
\newcommand{\TPPOTieRmsMwSd}{0.25}
\newcommand{\TPPOUnmetMax}{0.0}
\newcommand{\TPPOUnmetMean}{0.00}
\newcommand{\TPPOUnmetMwh}{0.00}
\newcommand{\TPPOUnmetMwhSd}{0.00}
\newcommand{\TPPOUvPuH}{0.90}
\newcommand{\TPPOUvPuHSd}{0.35}
\newcommand{\TPPOViolHours}{11.9}
\newcommand{\TPPOViolHoursSd}{5.5}
\newcommand{\TPPOVmaxDay}{1.024}
\newcommand{\TPPOVmaxDaySd}{0.022}
\newcommand{\TPPOVminDay}{0.925}
\newcommand{\TPPOVminDaySd}{0.006}
\newcommand{\TPRB}{1.9\times10^{-9}}
\newcommand{\TPUC}{1.9\times10^{-9}}
\newcommand{\TRBCost}{245.5}
\newcommand{\TRBCostSd}{31.6}
\newcommand{\TRBEViol}{0.000}
\newcommand{\TRBFRmsMhz}{11.48}
\newcommand{\TRBFRmsMhzSd}{1.16}
\newcommand{\TRBLossMwh}{3.416}
\newcommand{\TRBLossMwhSd}{0.261}
\newcommand{\TRBNadirMaxHz}{70.5}
\newcommand{\TRBNadirMaxHzSd}{7.6}
\newcommand{\TRBNadirMeanHz}{49.4}
\newcommand{\TRBNadirMeanHzSd}{6.0}
\newcommand{\TRBOvPuH}{0.001}
\newcommand{\TRBOvPuHSd}{0.005}
\newcommand{\TRBProj}{0}
\newcommand{\TRBProjEvents}{0.0}
\newcommand{\TRBProjEventsSd}{0.0}
\newcommand{\TRBPwExcess}{0.000}
\newcommand{\TRBRawClip}{0}
\newcommand{\TRBReserveMeanMw}{15.2}
\newcommand{\TRBReserveMeanMwSd}{1.4}
\newcommand{\TRBRocofMax}{0.137}
\newcommand{\TRBRocofMaxSd}{0.016}
\newcommand{\TRBSocDefPct}{13.58}
\newcommand{\TRBSocDefPctSd}{14.06}
\newcommand{\TRBThroughputMwh}{337}
\newcommand{\TRBThroughputMwhSd}{92}
\newcommand{\TRBTieRmsMw}{2.49}
\newcommand{\TRBTieRmsMwSd}{0.24}
\newcommand{\TRBUnmetMax}{0.0}
\newcommand{\TRBUnmetMean}{0.00}
\newcommand{\TRBUnmetMwh}{0.00}
\newcommand{\TRBUnmetMwhSd}{0.00}
\newcommand{\TRBUvPuH}{2.86}
\newcommand{\TRBUvPuHSd}{0.59}
\newcommand{\TRBViolHours}{12.5}
\newcommand{\TRBViolHoursSd}{2.1}
\newcommand{\TRBVmaxDay}{1.024}
\newcommand{\TRBVmaxDaySd}{0.016}
\newcommand{\TRBVminDay}{0.907}
\newcommand{\TRBVminDaySd}{0.007}
\newcommand{\TRedCostDAO}{11.0}
\newcommand{\TRedCostRB}{44.8}
\newcommand{\TRedCostUC}{71.6}
\newcommand{\TRedNadirMaxHzDAO}{1.3}
\newcommand{\TRedNadirMaxHzRB}{1.3}
\newcommand{\TRedNadirMaxHzUC}{12.2}
\newcommand{\TRedNadirMeanHzDAO}{5.0}
\newcommand{\TRedNadirMeanHzRB}{4.4}
\newcommand{\TRedNadirMeanHzUC}{13.2}
\newcommand{\TRedRocofMaxDAO}{3.2}
\newcommand{\TRedRocofMaxRB}{3.2}
\newcommand{\TRedRocofMaxUC}{14.1}
\newcommand{\TRedThroughputMwhDAO}{-28.2}
\newcommand{\TRedThroughputMwhRB}{-32.4}
\newcommand{\TRedThroughputMwhUC}{9.1}
\newcommand{\TRedUvPuHDAO}{13.5}
\newcommand{\TRedUvPuHRB}{68.4}
\newcommand{\TRedUvPuHUC}{87.8}
\newcommand{\TRedViolHoursDAO}{34.3}
\newcommand{\TRedViolHoursRB}{4.5}
\newcommand{\TRedViolHoursUC}{38.1}
\newcommand{\TUCCost}{477.2}
\newcommand{\TUCCostSd}{74.8}
\newcommand{\TUCEViol}{0.000}
\newcommand{\TUCFRmsMhz}{12.80}
\newcommand{\TUCFRmsMhzSd}{1.22}
\newcommand{\TUCLossMwh}{4.978}
\newcommand{\TUCLossMwhSd}{0.595}
\newcommand{\TUCNadirMaxHz}{79.2}
\newcommand{\TUCNadirMaxHzSd}{10.4}
\newcommand{\TUCNadirMeanHz}{54.4}
\newcommand{\TUCNadirMeanHzSd}{6.2}
\newcommand{\TUCOvPuH}{0.050}
\newcommand{\TUCOvPuHSd}{0.082}
\newcommand{\TUCProj}{833}
\newcommand{\TUCProjEvents}{27.8}
\newcommand{\TUCProjEventsSd}{5.0}
\newcommand{\TUCPwExcess}{0.000}
\newcommand{\TUCRawClip}{0}
\newcommand{\TUCReserveMeanMw}{0.0}
\newcommand{\TUCReserveMeanMwSd}{0.0}
\newcommand{\TUCRocofMax}{0.154}
\newcommand{\TUCRocofMaxSd}{0.024}
\newcommand{\TUCSocDefPct}{0.00}
\newcommand{\TUCSocDefPctSd}{0.00}
\newcommand{\TUCThroughputMwh}{491}
\newcommand{\TUCThroughputMwhSd}{19}
\newcommand{\TUCTieRmsMw}{2.51}
\newcommand{\TUCTieRmsMwSd}{0.23}
\newcommand{\TUCUnmetMax}{0.0}
\newcommand{\TUCUnmetMean}{0.00}
\newcommand{\TUCUnmetMwh}{0.00}
\newcommand{\TUCUnmetMwhSd}{0.00}
\newcommand{\TUCUvPuH}{7.40}
\newcommand{\TUCUvPuHSd}{1.49}
\newcommand{\TUCViolHours}{19.3}
\newcommand{\TUCViolHoursSd}{3.4}
\newcommand{\TUCVmaxDay}{1.044}
\newcommand{\TUCVmaxDaySd}{0.028}
\newcommand{\TUCVminDay}{0.835}
\newcommand{\TUCVminDaySd}{0.015}
\newcommand{\TWinsDAO}{30}
\newcommand{\TWinsRB}{30}
\newcommand{\TWinsUC}{30}
\newcommand{\TrainMin}{31}
\newcommand{\TrainSteps}{399\,360}
\newcommand{\TrkDAOBindHours}{2.8}
\newcommand{\TrkDAOGrossDec}{5.37}
\newcommand{\TrkDAOGrossDecRatio}{0.94}
\newcommand{\TrkDAOGrossFull}{5.70}
\newcommand{\TrkDAOGrossHour}{14.70}
\newcommand{\TrkDAOGrossHourRatio}{2.6}
\newcommand{\TrkDAOGrossWin}{0.25}
\newcommand{\TrkDAOMarginMin}{0.000}
\newcommand{\TrkDAONetDec}{-0.43}
\newcommand{\TrkDAONetFull}{-0.47}
\newcommand{\TrkDAONetHour}{-1.03}
\newcommand{\TrkDAONetWin}{-0.02}
\newcommand{\TrkDAOSatFrac}{0.00}
\newcommand{\TrkDAOSignFlipPct}{43}
\newcommand{\TrkPPOBindHours}{0.1}
\newcommand{\TrkPPOGrossDec}{8.29}
\newcommand{\TrkPPOGrossDecRatio}{0.94}
\newcommand{\TrkPPOGrossFull}{8.84}
\newcommand{\TrkPPOGrossHour}{25.44}
\newcommand{\TrkPPOGrossHourRatio}{2.9}
\newcommand{\TrkPPOGrossWin}{0.42}
\newcommand{\TrkPPOMarginMin}{0.000}
\newcommand{\TrkPPONetDec}{-0.69}
\newcommand{\TrkPPONetFull}{-0.73}
\newcommand{\TrkPPONetHour}{-1.67}
\newcommand{\TrkPPONetWin}{-0.03}
\newcommand{\TrkPPOSatFrac}{0.00}
\newcommand{\TrkPPOSignFlipPct}{39}
\newcommand{\TrkRBBindHours}{0.0}
\newcommand{\TrkRBGrossDec}{5.68}
\newcommand{\TrkRBGrossDecRatio}{0.94}
\newcommand{\TrkRBGrossFull}{6.02}
\newcommand{\TrkRBGrossHour}{15.98}
\newcommand{\TrkRBGrossHourRatio}{2.7}
\newcommand{\TrkRBGrossWin}{0.27}
\newcommand{\TrkRBMarginMin}{0.701}
\newcommand{\TrkRBNetDec}{-0.46}
\newcommand{\TrkRBNetFull}{-0.50}
\newcommand{\TrkRBNetHour}{-1.10}
\newcommand{\TrkRBNetWin}{-0.02}
\newcommand{\TrkRBSatFrac}{0.00}
\newcommand{\TrkRBSignFlipPct}{43}
\newcommand{\VOneDAOCost}{152.1}
\newcommand{\VOnePPOCost}{135.4}
\newcommand{\VOneRBCost}{245.0}
\newcommand{\VOneUCCost}{477.2}
\newcommand{\ValDAO}{156.1}
\newcommand{\ValDfSim}{-0.02381}
\newcommand{\ValDfTh}{-0.02381}
\newcommand{\ValLoss}{202.68}
\newcommand{\ValMode}{0.38}
\newcommand{\ValRB}{254.6}
\newcommand{\ValTieSim}{-5.000}
\newcommand{\ValTieTh}{-5.000}
\newcommand{\ValUC}{496.9}
\newcommand{\ValVerr}{1.2\times10^{-10}}
\newcommand{\ValVmin}{0.9131}

\newcommand{\rc}[1]{\par\noindent\textbf{Comment #1}\par}
\newcommand{\resp}{\par\noindent\textbf{Response:}\ }
\newcommand{\chg}{\par\noindent\textbf{Changes in the manuscript:}\ }
\newcommand{\quoted}[1]{\begin{quote}\itshape #1\end{quote}}
\begin{document}

\begin{center}
{\large\bfseries Response to the Reviewers}\\[4pt]
Resubmission of the manuscript previously titled ``Coordinated V2G Voltage and Frequency-State Regulation in a Two-Area System: A Gaussian Actor--Critic Simulation Framework'' (manuscript ID: \textcolor{red}{[previous ID]})\\[2pt]
New title: ``Two-Time-Scale Coordination of Vehicle-to-Grid Voltage Support and Frequency Reserve in a Two-Area System Using Proximal Policy Optimization''
\end{center}

Dear Editor and Reviewers,

We thank both reviewers for their careful and direct assessment. They
identified the central weakness of the previous version correctly: the
controller was never trained, the network and frequency models were
algebraic surrogates, and the comparison curves were derived rather
than simulated. We agree with this assessment. Rather than revising
the text, we rebuilt the study from the model upward. The main changes
are:

\begin{enumerate}[leftmargin=1.5em]
\item The surrogate voltage equations were replaced by an AC power flow
on the IEEE 33-bus feeder (Baran--Wu data) in each area, solved by a
backward/forward sweep and verified against pandapower (maximum voltage
difference $\ValVerr$~p.u.).
\item The algebraic frequency equation was replaced by a dynamic
two-area load-frequency control model with inertia, damping,
governor--turbine dynamics, a tie line with synchronizing coefficient,
ACE-based AGC, and a saturated V2G droop loop, integrated at a 20-ms
step. Hourly scheduling and second-level frequency dynamics are now
separate time scales.
\item All actions and states have physical units: fleet size, battery
capacity, charger rating, efficiencies, availability, driving energy,
area base, and tie-line dynamics are specified (Table~\ref{tab:params} of the
manuscript). Sampled actions are clipped before use, and set-points are
projected onto charger and energy limits; raw-action clippings and
physical projections are counted separately.
\item PPO is now actually trained (Stable-Baselines3) with \NSeed{}
independent seeds for \TrainSteps{} steps each, with a fixed validation
set and a stated convergence criterion. All diagnostic or analytical
curves of the previous version (old Figs.~3 to 8 and 13) were removed.
\item Three benchmark controllers are simulated independently under
identical scenarios, disturbances, network, and constraints:
uncontrolled charging, a conventional rule-based volt-watt and droop
controller, and a day-ahead optimization with perfect foresight.
\item Results are reported on 30 held-out days and 20
out-of-distribution stress days with quantitative metrics, standard
deviations across days and seeds, and paired significance tests. A
reward ablation is included.
\item The energy exchanged by the frequency droop loop is included in
the fleet energy balance with both efficiencies. Four treatments of
this term were compared against exact 3600-s simulations of every
evaluated hour; the adopted treatment tracks the exact value to within 6\% and the
results are insensitive to the choice (Section~\ref{sec:regenergy}).
\item Raw-action saturation of the Gaussian policy is reported per
action dimension and compared with a bounded (tanh-squashed, gSDE)
policy trained at the same budget (Section~\ref{sec:clip}).
\item All policies were trained from scratch in the final model
version; every number in the manuscript is recomputed from the raw
result files by an independent verification script, and the code,
trained policies, data, checksums, and exact package versions are
released.
\end{enumerate}

Because the model, the method, and all results changed, we are
submitting the work as a new manuscript and summarize below how each
comment was addressed. Section, table, and figure numbers refer to the
new manuscript.

\section*{Reviewer 1}

\rc{1.1} \emph{Implement and document a complete training procedure,
with genuine training and held-out evaluation results, or redefine the
contribution.}
\resp We implemented the complete PPO procedure (clipped surrogate
objective, generalized advantage estimation, minibatch updates) with
Stable-Baselines3. Training uses \NSeed{} seeds, four parallel
environments, and randomized daily scenarios. Every 6144 steps the
deterministic policy is evaluated on ten fixed validation days, and
the policy is finally tested on 30 held-out days and 20 stress days
that are never used for training or model selection. The final
validation cost changes by at most \ConvChangePct\% over the last
50\,000 steps.
\chg Section~\ref{sec:ppo} (training procedure), Table~\ref{tab:params} (hyperparameters),
Fig.~\ref{fig:learning} (learning curves on training and validation days), Section~\ref{sec:training}.

\rc{1.2} \emph{Implement a documented network model with branch
parameters, nodal injections, and a reproducible power-flow solution,
or stop presenting the surrogate as IEEE 33-bus.}
\resp Each area now contains the IEEE 33-bus feeder with the original
branch impedances and nominal loads (12.66~kV), taken from the
pandapower case library. Net nodal injections combine load, PV, and V2G
power at specified buses, and the power flow is solved by the direct
BIBC/BCBV method. The base case reproduces the published values
(minimum voltage \ValVmin~p.u.\ at bus 18, loss \ValLoss~kW), and the
solver agrees with pandapower Newton--Raphson to within $\ValVerr$~p.u.
over 20 random operating points.
\chg Section~\ref{sec:feeders}, Eqs.~(\ref{eq:netload}) and (\ref{eq:bfs}), Fig.~\ref{fig:feeder}, Section~\ref{sec:verif}.

\rc{1.3} \emph{Provide a physically justified frequency model, define
the units of the coefficient 0.15, and relate actions to power
imbalance.}
\resp The coefficient 0.15 and the algebraic frequency equation were
removed. The new model (Eqs.~(\ref{eq:swing}) to (\ref{eq:evdroop})) is the standard two-area LFC
model in per unit of a 500-MW area base and 50~Hz, with inertia
constants that decrease with PV output, load damping, governor and
turbine time constants, droop, a tie line with synchronizing
coefficient $T_{12}$, ACE, and integral AGC. The V2G action enters
physically: the reserve in MW sets the droop gain and the saturation
of the fleet's response. The model reproduces the closed-form
steady-state deviation ($\ValDfSim$~Hz simulated against $\ValDfTh$~Hz
theoretical) and tie-line change.
\chg Section~\ref{sec:lfcmodel}, Table~\ref{tab:params}, Section~\ref{sec:verif}.

\rc{1.4} \emph{Simulate the controlled and baseline cases independently
under identical conditions, include an uncontrolled case and a
conventional controller, with quantitative indices.}
\resp All controllers now run in the same environment with the same
scenarios, the same per-hour disturbance seeds, the same network, and
the same constraints. The benchmarks are uncontrolled charging (UC), a
rule-based volt-watt and droop controller (RB), and a day-ahead
optimization with perfect foresight (DAO). Sixteen quantitative indices
are reported with standard deviations, and paired Wilcoxon tests are
given. On the test days, PPO reduces the daily cost by \TRedCostDAO\%,
\TRedCostRB\%, and \TRedCostUC\% relative to DAO, RB, and UC ($p=\TPDAO$).
\chg Sections~\ref{sec:bench} and \ref{sec:protocol}, Tables~\ref{tab:test} and \ref{tab:stress}, Figs.~\ref{fig:day} to \ref{fig:box}.

\rc{1.5} \emph{The active-power-loss comparison in Fig.~13 cannot be
reproduced.}
\resp The figure has been removed. Feeder losses are now computed from
the branch currents of the power-flow solution and are reported as one
of the indices in Table~\ref{tab:test}, for every controller.
\chg Table~\ref{tab:test}.

\rc{1.6} \emph{Define the physical units and scaling of each action,
impose explicit bounds, and demonstrate compliance with fleet power,
converter, and inter-area limits.}
\resp The action now consists of a normalized set-point and a reserve
share per area. Eqs.~(\ref{eq:decode}) and (\ref{eq:reserve}) map them to MW using the plugged-in
charger rating, and the set-point is projected onto the charger and
energy limits of Eqs.~(\ref{eq:pdis}) and (\ref{eq:pch}); the reserve is limited so that the
droop deployment also respects these limits. Actions sampled from the
Gaussian policy are clipped to $[-1,1]$ before this mapping. Raw-action
clipping and physical projection are counted separately: the PPO policy
mean lies outside $[-1,1]$ for \RawClipFracPct\% of action components
(saturation at the design limits), while the physical projection is
never active for PPO. The combined scheduled and droop power is audited
at every 20-ms step and never exceeds the charger or energy limits
(Section~\ref{sec:regenergy}, Table~\ref{tab:regenergy}). Tie-line power is a dynamic state of the LFC
model, and its RMS deviation is reported. The ``MW'' labels
of the old figures, which had no physical basis, are gone.
\chg Section~\ref{sec:obsact}, Section~\ref{sec:regenergy}, Tables~\ref{tab:test} and \ref{tab:regenergy}.

\rc{1.7} \emph{Provide an energy-consistent fleet model.}
\resp The fleet model (Eqs.~(\ref{eq:prate}) to (\ref{eq:regenergy})) specifies 300 vehicles per
feeder, 60-kWh batteries, 7.4-kW chargers, charging and discharging
efficiencies of 0.92, a 1-h step, SoC limits of 0.20 and 0.90,
time-varying plug-in availability, and 9~kWh of daily driving energy
per vehicle. SoC-dependent power limits follow from the remaining
energy. The energy exchanged by the frequency droop loop is computed
from the fast-layer simulations of every hour, with both efficiencies,
and added to the energy balance (Eq.~(\ref{eq:energy})); four
treatments of this term are compared with exact 3600-s simulations in
Section~\ref{sec:regenergy} and its size is reported in
Table~\ref{tab:regenergy}. Raw-action clippings, projection events, and unmet mobility energy are recorded and
reported.
\chg Section~\ref{sec:fleet}, Tables~\ref{tab:params}, \ref{tab:test}, and \ref{tab:regenergy}.

\rc{1.8} \emph{Justify the penalty weights, specify the voltage and
frequency limits, penalize both lower and upper voltage deviations,
and demonstrate the effect of the reward design.}
\resp The reward now penalizes both undervoltage and overvoltage
outside 0.95 to 1.05~p.u., the frequency nadir, the tie-line deviation,
battery throughput, unmet mobility energy, and the end-of-day SoC
deficit (Eq.~(\ref{eq:cost})). The weights were set so that voltage and frequency
costs are of the same order on an uncontrolled day, which is stated in
the text. A reward ablation (Table~\ref{tab:ablation}) retrains the policy without the
overvoltage term, without the frequency terms, and without the
throughput term, and evaluates each with the full cost against a
reference trained with the same budget. Removing the frequency terms
clearly reduces the reserve (\AblNoFReserveMeanMw~MW instead of
\AblRefReserveMeanMw~MW) and deepens the nadir, which confirms
that the learned reserve allocation comes from these terms. Removing
the overvoltage or throughput term gives no difference beyond the
seed-to-seed spread; the lower mean cost of these variants is traced,
through a cost decomposition, to the slowly learned terminal-SoC term
and the shorter training budget, not to the removed term. Halving and
doubling $w_V$ and $w_f$ move reserve, undervoltage, and terminal SoC
in the expected directions (Table~\ref{tab:sensitivity}).
\chg Section~\ref{sec:reward}, Section~\ref{sec:ablation}, Tables~\ref{tab:ablation} and \ref{tab:sensitivity}.

\rc{1.9} \emph{Report genuine training episodes, learning settings,
convergence criteria, and evaluation on unseen conditions with multiple
seeds and varying load, solar, and fleet participation.}
\resp Done as described under Comment 1.1. The scenario generator
varies load level, cloud cover, intra-day PV variability, fleet
participation, initial SoC, and disturbance size. Test days come from
an independent random stream, and the stress days use heavier load, lower
participation, and larger disturbances than any training day, combined
with clear-sky PV at the upper end of the training range. Results are given as mean $\pm$ standard deviation across
days, with the seed-to-seed standard deviation in parentheses.
\chg Sections~\ref{sec:profiles}, \ref{sec:ppo}, \ref{sec:protocol}, and \ref{sec:training} to \ref{sec:stress}.

\rc{1.10} \emph{Identify the specific technical innovation and compare
with prior work.}
\resp The revised contribution is the joint, learned allocation of
each fleet's converter capability between hourly voltage support and
frequency reserve, evaluated on a branch-resolved distribution model
and a dynamic two-area LFC model at the same time. Voltage-oriented
V2G studies do not model frequency dynamics, and frequency-oriented
V2G studies do not evaluate feeder voltages; aggregate PEV models for
primary frequency control fix the participation in advance. The
results show the value of this allocation: PPO holds more reserve than
the fixed-share benchmarks where vehicles are connected and inertia is
low, releases it when voltage support needs the converter, and still
achieves lower undervoltage than a perfect-foresight schedule.
Where the gains are small, as for the frequency nadir (5\% relative to
the day-ahead optimization), the manuscript says so.
\chg Section~\ref{sec:introduction} (rewritten), contribution list, Section~\ref{sec:alloc}.

\textbf{Additional question 4 (references).} We added references on
aggregate PEV frequency control models, two-time-scale RL voltage
control, low-inertia systems, LFC modeling, the power-flow method,
pandapower, Stable-Baselines3, generalized advantage estimation, and
reproducibility in deep RL.

\section*{Reviewer 2}

\rc{2.1} \emph{Train the controller, report real learning curves from
several seeds, and remove Figs.~3 to 8.}
\resp PPO is trained with \NSeed{} seeds; Fig.~\ref{fig:learning} shows the real
learning curves on training and validation episodes. Old Figs.~3 to 8
were removed.

\rc{2.2} \emph{Explain how Fig.~12 was produced; under Eqs.~(3) to (6)
the minimum voltage cannot reach 0.955~p.u.; conflict with Fig.~14.}
\resp The reviewer is right. The previous voltage equations and figures
were inconsistent and have been withdrawn. All voltages in the new
manuscript come from the AC power flow, and every voltage figure and
table is generated from the same simulation records by the released
scripts.

\rc{2.3} \emph{Explain the gap between Figs.~10 and 11 (SoC change
inconsistent with Eq.~(9)); exact zeros in Fig.~10; noisy Fig.~16.}
\resp These figures were inconsistent with the previous model and have
been withdrawn. In the new version, SoC follows the energy balance of
Eq.~(\ref{eq:energy}), which can be checked against the reported set-points
(Fig.~\ref{fig:day}(c) and (e)).

\rc{2.4} \emph{Use a real network or rename the environment.}
\resp See the response to Comment 1.2.

\rc{2.5} \emph{Use a real two-area frequency model with inertia,
damping, governor, per-area frequency, tie-line power with a
synchronizing coefficient, and ACE, at a suitable time step; separate
the time scales.}
\resp Implemented exactly as suggested (Section~\ref{sec:lfcmodel}): per-area
frequency, inertia, damping, governor and turbine, tie line with
$T_{12}$, ACE with integral AGC, a 20-ms RK4 step within a 60-s window
each hour, and an hourly scheduling layer for voltage and energy.

\rc{2.6} \emph{Give the actions physical units; clip or squash actions
before env.step(); add an overvoltage penalty; consider cost and
degradation terms.}
\resp See Comments 1.6 to 1.8. Actions are clipped before use, an
overvoltage penalty is included, and a throughput-based degradation
proxy is part of the reward (its effect is shown in the ablation). We
also compared clipping with squashing: a tanh-squashed Gaussian policy
with generalized state-dependent exploration was trained with three
seeds at the ablation budget. It removes saturation by construction
but reaches a higher cost than the clipped policy at that budget
(Table~\ref{tab:clip}); the clipped policy is therefore retained and
its saturation per action dimension is reported. An electricity-cost
term was not added, because the paper is about the technical
allocation of converter capability; it is listed as future work.

\rc{2.7} \emph{Use a realistic daily load profile; the $\sin^2$ shape
peaks at midnight and noon.}
\resp Correct; the previous profile had a 12-h period. The new profiles
(Fig.~\ref{fig:profiles}) have a residential evening peak at 19:00 to 21:00 and a mixed
commercial and residential daytime plateau. They are synthetic and
are randomized in level and noise for each scenario.

\rc{2.8} \emph{Report voltage violation hours and magnitude, frequency
deviation statistics, SoC at end of day, and energy cycled across seeds
and scenarios.}
\resp All of these are reported in Table~\ref{tab:test} for the test days (Table~\ref{tab:stress}
gives the main subset for the stress days, and the full stress-day
records are in the released data): violation hours,
integrated undervoltage and overvoltage, minimum and maximum voltage,
mean and worst nadir, RMS frequency deviation, RoCoF, tie-line
deviation, terminal SoC deficit, fleet throughput, unmet mobility energy, projection events, and
losses, across \NSeed{} seeds and 30 test or 20 stress days.

\rc{2.9} \emph{Share the code and data.}
\resp The complete code (environment, models, verification, training,
evaluation, figure scripts), the trained policies, and the result files
are available at \url{https://github.com/kenoto99/v2g-ppo-two-area}
(release v3.0.1) and archived on Zenodo (DOI 10.5281/zenodo.23851585),
as given in the Data and Code Availability statement.

\rc{2.10} \emph{Cut the repeated disclaimers once real results are in;
fix the keywords and subject category.}
\resp The disclaimers were removed because the underlying limitations
were removed. Remaining limitations are stated once, in Section~\ref{sec:limits}.
The keywords now are: electric vehicles, frequency regulation,
load-frequency control, proximal policy optimization, reinforcement
learning, vehicle-to-grid, voltage regulation.

\textbf{Additional question 1.} The previous evaluation was a single
seeded rollout of an untrained policy. The new results are from
trained policies evaluated on unseen days.


\section*{Additional changes made after a technical review of the revised manuscript}

Before resubmission the revised manuscript was examined by an
independent technical reviewer, whose eight comments are addressed
below. The analyses added in response reuse the trained policies and
the test-day scenarios; no training run was repeated. The author list
was also updated with the agreement of all authors: Wasan Phetphimoon
is no longer an author, and Kritphon Phanrattanachai (Phetchabun
Rajabhat University) was added as the sixth author.

\rc{A.1} \emph{Fairness of the comparison with DAO: PPO adjusts both
the V2G power and the reserve share, whereas DAO fixes the share at
0.3 and does not optimize the frequency part.}

\resp We agree that the fixed share could have favored PPO. The
benchmark was therefore re-optimized for every share
$\rho\in\{0,0.1,\dots,0.6\}$ on the same test and stress days, and a
per-day oracle (the share with the lowest realized cost on each day,
an upper bound on any tuned fixed share) was added. The fixed share
used before was close to, but not at, the best fixed share
($\rho=\DaoRhoBest$); PPO remains \RedCostDaoRhoBest\% below the best
fixed share and \RedCostDaoOracle\% below the oracle on the test days
(lower on \WinsDaoOracle{} of 30 days, $p=\PDaoOracle$) and
\RedCostDaoOracleS\% below it on the stress days. The sweep also shows
the trade-off that no fixed share resolves (nadir improves and
undervoltage worsens with $\rho$), which PPO avoids by moving the share
hour by hour. The structural difference between the two controllers
(DAO's objective excludes the frequency terms) is now stated
explicitly.

\chg Section~\ref{sec:bench} (last paragraph); new
Section~\ref{sec:daorho} and Table~\ref{tab:daorho}; Conclusion.

\rc{A.2} \emph{The daily cost of Table~\ref{tab:test} is a weighted sum
of technical penalties, not an economic cost; its meaning and units
should be stated before relative reductions are interpreted.}

\resp The quantity is now defined as the daily performance cost, a
dimensionless weighted penalty with no price, reserve payment, or
degradation cost, in Section~\ref{sec:reward}; the table row, the
learning-curve axis, the first sentence of
Section~\ref{sec:testres}, and the abstract say so, and the relative
reductions are described as relative technical penalty under the
chosen weights.

\chg Section~\ref{sec:reward} (new last paragraph); abstract;
Section~\ref{sec:testres}; row label of Tables~\ref{tab:test} and
\ref{tab:stress}; Fig.~\ref{fig:learning}.

\rc{A.3} \emph{The voltage results do not meet the band in every hour
(PPO minimum voltage \TPPOVminDay~p.u., \TPPOViolHours{} violation hours
per day); this should be presented as a reduction of severity, not a
guarantee.}

\resp Agreed. The voltage paragraph of Section~\ref{sec:testres} now
states that the results are a reduction in depth and duration of
undervoltage and not compliance in every hour, gives the number of
area-hours that remain outside the band, and explains why the feeder
itself sets the floor (nominal Baran--Wu minimum of \ValVmin~p.u.\
with no EVs; the smallest plugged-in fleet during the evening peak).
The abstract and the conclusion carry the same qualification.

\chg Section~\ref{sec:testres}, paragraph \emph{Voltage}; abstract;
Conclusion.

\rc{A.4} \emph{Terminal SoC deficit (\TPPOSocDefPct\% on test days,
\SPPOSocDefPct\% on stress days): explain the consequence for driving
energy and consider a hard terminal constraint or consecutive-day
tests.}

\resp The deficit is now translated into energy per vehicle
(\TPPOSocDefKwhVeh~kWh, about \TPPOSocDefKm~km or \TPPOSocDefDrivePct\%
of a day's driving energy on the test days), and the unmet mobility
energy of the day is reported alongside. More importantly, the 30 test
days were run consecutively with the terminal energy carried over
(new Table~\ref{tab:chained}). Costs rise for all controllers, and the
residential fleet drifts to a lower SoC over the month with all three;
PPO degrades least (unmet mobility energy \ChPPOUnmetMwhTotal~MWh on
\ChPPOUnmetDays{} days, against \ChDAOUnmetMwhTotal~MWh on
\ChDAOUnmetDays{} days for DAO), but the drift shows that the soft
terminal penalty does not guarantee a sustainable energy state. We
state this plainly, identify the hard terminal constraint (constrained
RL or a projection at the last plugged-in hours) as the remedy, and
did not implement it in this revision.

\chg Section~\ref{sec:testres}, paragraph \emph{Energy and fleet};
new Table~\ref{tab:chained}; Conclusion; Supplementary S6 and S7.

\rc{A.5} \emph{Limits of the aggregated EV model: how are parked and
driving vehicles separated, and how does the single-SoC assumption
affect deliverable power?}

\resp Section~\ref{sec:fleet} now explains that the aggregation
separates parked and driving vehicles in power (through the plugged-in
fraction) but not in energy (driving energy is withdrawn from the
common store while vehicles are away), that all plugged-in vehicles
are assumed to share the fleet-average SoC, and that the
energy-feasible limits are exact for the aggregate but optimistic
for the deliverable power of a dispersed fleet near its SoC limits,
most of all at the end of the evening discharge. A vehicle-level
allocation layer is listed as future work.

\chg Section~\ref{sec:fleet}; Section~\ref{sec:limits} item (iii);
Supplementary S7.

\rc{A.6} \emph{Coupling between the frequency reserve and the feeder
power flow: when reserve is activated, do feeder voltages and currents
stay within limits?}

\resp This was not checked before and the check changes the picture.
Every test-day hour was re-solved at $P^{\mathrm{set}}\pm R$ and at the
realized peak droop power (new Section~\ref{sec:activation} and
Table~\ref{tab:activation}). Upward activation helps the voltage;
full downward activation (the fleet absorbing the whole reserve)
would take \ActPPONewViolFracDown\% of the PPO area-hours that are
inside the band at the schedule out of it, with a minimum voltage of
\ActPPOVminDown~p.u. The simulated disturbances never reach that
level: the realized peak droop power averages \ActPPOPeakActFracMean\%
of the committed reserve, at which \ActPPONewViolFracDownreal\% of the
area-hours leave the band for a few seconds, with an additional drop
of at most \ActPPODvminDownrealMax~p.u. Branch currents change little
and the case data carry no thermal rating. The exposure is small for
the modeled disturbance class but real and growing with the committed
reserve, so we report it as a limitation and identify a voltage-aware
(asymmetric) reserve rule as the remedy.

\chg New Section~\ref{sec:activation} and Table~\ref{tab:activation};
Conclusion; Supplementary S6.

\rc{A.7} \emph{Effect of raw-action saturation (\ClipDimAll\%) on
exploration, gradients, and transfer; limits of the squashed-policy
comparison.}

\resp Section~\ref{sec:clip} now discusses the three effects
separately (exploration collapses to the boundary action, the
likelihood-ratio update keeps pushing the mean outward with no
corrective signal from the environment, and a saturated output is
benign only while the deployed limits coincide with the training
limits) and states the limits of the squashed-policy comparison
(three seeds, half budget, different exploration mechanism, no
re-tuning, one parameterization), so that it is read as "the bounded
variant we tried did not help under equal conditions" rather than as
a general statement.

\chg Section~\ref{sec:clip}, two paragraphs.

\rc{A.8} \emph{Realism and scope: discuss forecast errors,
communication delay, charger heterogeneity, converter dynamics, and
parameter uncertainty.}

\resp The Limitations section now has a dedicated paragraph with
these five items, each with its consequence for the reported results
and the step needed before field use (training on forecast errors,
delays, and parameter ranges; a vehicle-level dispatch layer; a fuller
converter model). The stress days are explicitly described as varying
load, participation, and disturbance level but not forecast accuracy.

\chg Section~\ref{sec:limits} (new paragraph, items (i) to (v));
Conclusion; Supplementary S7.

We thank the reviewers again. Their comments led to a substantially
stronger study.

Sincerely,\\
Krischonme Bhumkittipich, on behalf of all authors

\end{document}
